\documentclass[aspectratio=169,10pt,spanish]{beamer}
\input{../../../_preambulo_beamer.tex}
\input{../../../_comandos_beamer.tex}
\selectlanguage{spanish}
\shorthandoff{"}
\title{L06: Técnicas de conteo I}
\subtitle{Ruta de clase -- 40 minutos}
\author{}
\date{}
\begin{document}

\begin{frame}{Técnicas de conteo I}
\framesubtitle{Ruta de clase -- 40 minutos}
\textbf{Pregunta:} ¿Qué cuenta como un resultado y cómo contamos sin duplicarlo?
\begin{block}{Al terminar podrás}
Aplicar la regla del producto, distinguir permutaciones de combinaciones, comprobar una fórmula por enumeración y rechazar un sobreconteo.
\end{block}
\smallskip
Todas las instalaciones, personas, funciones y restricciones son supuestos didácticos sintéticos.
\end{frame}

\begin{frame}{Ruta de 40 minutos}
\small
\begin{tabular}{p{.42\linewidth}rp{.39\linewidth}}
\toprule Tema & Min & Comprobación observable \\ \midrule
Definir un resultado & 3 & Decidir si importan los roles.\\
Regla del producto & 7 & Contar 24 configuraciones.\\
Permutaciones & 8 & Contar 120 asignaciones.\\
Combinaciones & 7 & Contar 20 comités sin roles.\\
Restricciones & 11 & Rechazar 420 y derivar 99.\\
Integración & 4 & Elegir y justificar el método.\\ \bottomrule
\end{tabular}
\par\medskip
Ejecuta las celdas preparadas de \texttt{leccion\_06\_compacta.ipynb} en cada paso. El cómputo tarda segundos; el tiempo incluye razonamiento y comprobaciones. El estudio extenso permanece en la libreta y el mazo completos.
\end{frame}

\begin{frame}{Primero define el resultado}
\textbf{Experimento 1:} Configurar una revisión sintética eligiendo un centro de cumplimiento, una etapa del proceso y un tipo de evidencia.
\[
\omega=(\text{centro},\text{etapa},\text{tipo de evidencia}).
\]
Hay tres centros, cuatro etapas y dos tipos de evidencia. Dos ternas que difieren en alguna posición son configuraciones distintas.
\par\medskip
\textbf{Pregunta:} Si asignamos personas a funciones con nombre, ¿intercambiar a dos personas cambia el resultado?
\par\smallskip
\textbf{Comprobación:} Sí. Las funciones distinguen los resultados; por ello usaremos permutaciones.
\end{frame}

\begin{frame}[fragile]{Regla del producto: elecciones sucesivas}
\begin{columns}[T]
\begin{column}{.47\textwidth}
\small
Una elección de cada etapa produce
\[
N=\prod_i n_i=3\times4\times2=24.
\]
\begin{lstlisting}
from itertools import product
espacio = list(product(centros,
                       etapas, evidencias))
assert len(espacio) == 24
\end{lstlisting}
La libreta verifica ocho configuraciones por centro. Los factores cuentan componentes de una terna, no casos alternativos.
\end{column}
\begin{column}{.51\textwidth}
\centering
\includegraphics[width=\linewidth]{../figuras/conteo_01_regla_producto.png}
\end{column}
\end{columns}
\end{frame}

\begin{frame}{Permutaciones: funciones distintas}
\begin{columns}[T]
\begin{column}{.46\textwidth}
\small
Asignamos las funciones Líder, Validación y Reportes a tres de seis analistas. Las posiciones son distintas; importa el orden:
\[
P(6,3)=\frac{6!}{(6-3)!}=6\times5\times4=120.
\]
La libreta enumera 120 asignaciones ordenadas. Intercambiar Líder y Reportes crea otro resultado.
\par\smallskip
\textbf{Comprobación:} ¿Dividimos aquí entre $3!$? No: las funciones distinguen las asignaciones.
\end{column}
\begin{column}{.52\textwidth}
\centering
\includegraphics[width=\linewidth]{../figuras/conteo_02_permutaciones.png}
\end{column}
\end{columns}
\end{frame}

\begin{frame}{Combinaciones: un comité sin roles}
\begin{columns}[T]
\begin{column}{.48\textwidth}
\small
Elegimos tres de los mismos seis analistas para un comité sin funciones internas. El orden ya no crea un resultado nuevo:
\[
\binom{6}{3}=\frac{6!}{3!3!}=20.
\]
Las 120 asignaciones ordenadas se reducen a 20 comités porque cada uno aparece en $3!=6$ órdenes. La libreta enumera los 20 subconjuntos.
\par\smallskip
\textbf{Prueba de decisión:} ¿El resultado contiene posiciones con nombre o sólo un conjunto de personas?
\end{column}
\begin{column}{.50\textwidth}
\centering
\includegraphics[width=\linewidth]{../figuras/conteo_03_combinaciones.png}
\end{column}
\end{columns}
\end{frame}

\begin{frame}{Restricciones: primero define el universo}
\small
Elegimos un grupo de cuatro personas entre cinco analistas de operaciones $O_1,\ldots,O_5$ y cuatro especialistas de datos $D_1,\ldots,D_4$.
\begin{enumerate}
\item El grupo debe incluir ambas disciplinas.
\item $O_1$ y $D_1$ no pueden estar juntos en el grupo.
\end{enumerate}
No hay funciones internas en el grupo, así que el universo contiene
\[
|U|=\binom{9}{4}=126
\]
grupos sin orden. Una cuenta mayor que 126 es imposible.
\par\smallskip
\textbf{Comprobación:} ¿Qué método anterior cuenta el universo sin restricciones? Combinaciones.
\end{frame}

\begin{frame}{Por qué falla la tentación de 420}
\small
La propuesta $5\times4\times\binom{7}{2}=420$ primero designa a un analista de operaciones y a uno de datos, y luego elige dos personas más.
\begin{alertblock}{Comprobación de plausibilidad}
$420>126=|U|$, por lo que no puede contar grupos distintos de cuatro personas.
\end{alertblock}
Las personas designadas no tienen funciones especiales en el grupo final. Un grupo con dos personas de cada disciplina puede generarse $2\times2=4$ veces.
\par\smallskip
\textbf{Corrección:} Contar cada grupo una sola vez y restar los inválidos.
\end{frame}

\begin{frame}{El complemento produce 99}
\begin{columns}[T]
\begin{column}{.49\textwidth}
\small
Grupos de una sola disciplina: $\binom54+\binom44=5+1=6$.
\[
\text{Mixtos}=126-6=120.
\]
Grupos que contienen a $O_1$ y $D_1$: elegimos otras dos personas de las siete restantes, $\binom72=21$. Todos esos grupos ya son mixtos.
\[
\text{Válidos}=120-21=99.
\]
La libreta examina los 126 subconjuntos candidatos y confirma 99 válidos.
\end{column}
\begin{column}{.49\textwidth}
\centering
\includegraphics[width=\linewidth]{../figuras/conteo_04_restricciones.png}
\end{column}
\end{columns}
\end{frame}

\begin{frame}{Elige el método y luego compruébalo}
\small
\begin{enumerate}
\item Una elección por cada etapa sucesiva: \textbf{regla del producto}.
\item Personas en funciones distintas: \textbf{permutaciones}.
\item Comité sin roles internos: \textbf{combinaciones}.
\item Restricciones que interactúan: definir el universo, contar grupos inválidos y \textbf{restar sin errores de traslape}.
\end{enumerate}
\par\medskip
\textbf{Comprobación:} ¿Por qué la enumeración respalda 99 pero no elige la fórmula por nosotros? Compara los resultados generados con el modelo; primero debemos definir un resultado y las restricciones.
\par\smallskip
\textbf{Siguiente:} L07 cuenta categorías repetidas y pregunta cuándo un conteo no determina una probabilidad.
\end{frame}
\end{document}
